\section{Introduction}

Control Barrier Functions (CBFs) provide a principled approach to safety-critical control by certifying that a system remains within a safe set defined by a scalar function $h(x) \geq 0$~\citep{ames2019control}. Neural networks have emerged as powerful function approximators for CBFs due to their expressiveness, but verifying that a neural CBF (NCBF) actually satisfies the barrier condition $\dot{h}(x) + \alpha(h(x)) \geq 0$ for all states is challenging due to the nonlinear activation functions.

Prior work on NCBF verification employs linear relaxation techniques such as CROWN~\citep{ruppert1998ball, hwang2020crown} or LBP~\citep{liu2021 certify} to compute rigorous bounds on the neural network over a region. However, neural networks trained to approximate CBFs often fail verification out of the box, and na\"ive fine-tuning fails because gradient descent may violate already-verified safe regions.

In this paper, we propose a \textbf{progressive depth phased repair} algorithm that addresses this challenge through two key innovations:

\begin{enumerate}
    \item We classify verification failures into \emph{definitive violations} ($F_h$: $h(x) > 0$ in unsafe, $F_{safe}$: CBF condition violated) and \emph{uncertain regions} ($F_{depth}$: depth limit reached, $F_{split}$: cannot split). This distinction is crucial because uncertain regions require deeper verification to resolve.
    
    \item We employ a two-phase repair strategy: Phase 1 fixes definitive violations at shallow depth where violations are clearly detectable, then Phase 2 progressively increases verification depth to address uncertain regions that become verifiable at deeper levels.
\end{enumerate}

To enable efficient gradient-based repair without backpropagating through nondifferentiable activations, we use a \emph{randomized smoothing Jacobian} estimator that computes gradient information through the reparameterization trick. We then project gradient updates using quadratic programming to ensure that already-verified safe regions remain verified after the update.

Our key theoretical contributions are:
\begin{itemize}
    \item A formulation of the NCBF repair problem as constrained optimization with QP projection
    \item An analysis of the randomized smoothing Jacobian estimator for CBF bounds, including unbiasedness and variance bounds
    \item A progressive depth strategy with formal stopping criteria based on plateau detection
\end{itemize}

Experiments on multiple dynamical systems (Barrier1, Barrier2, Barrier3, Simple2D) with different activation functions (ReLU, Tanh, Sigmoid) demonstrate that our method achieves higher certified percentages than prior repair approaches.

\section{Related Work}

\subsection{Neural Network Verification}

Neural network verification has seen significant progress with linear relaxation-based methods. CROWN~\citep{ruppert1998ball, hwang2020crown} computes bounds by linearizing each layer and propagating bounds forward. LBP~\citep{liu2021 certify} uses bound propagation with layer-wise abstractions. These methods provide sound but incomplete verification---they may report ``unknown'' for regions where verification is inconclusive.

Semidefinite programming (SDP) relaxation provides an alternative approach to verification~\citep{fazel2018global}, but scales poorly to large networks. Our method builds on CROWN linearization but focuses on repair rather than just verification.

\subsection{Neural Network Repair}

Neural network repair addresses cases where verification fails. Prior work includes:
\begin{itemize}
    \item \citep{liu2021repair} who propose repair methods that guarantee satisfaction of desired properties
    \item \citep{islam2024repair} who use gradient-based optimization with constraint projection
    \item \citep{dutt2023consistent} who ensure consistency between original and repaired networks
\end{itemize}

Our work differs in that we specifically target the NCBF setting where the repair must preserve verified safe regions while correcting violations.

\subsection{Control Barrier Functions}

CBFs were introduced in~\citep{ames2019control} and have been extended to learning-based settings in~\citep{jwang2020learning}. Neural CBFs are studied in~\citep{cheng2019end, agarwal2019learning} where the barrier function is represented as a neural network. The challenge of verifying neural CBFs is addressed in~\citep{srinivasan2020verification, chang2019verified}.

To the best of our knowledge, our work is the first to propose a progressive depth phased repair strategy specifically designed for NCBF verification, where the key insight is that verification depth directly affects which violations are detectable.